\chapter{Constant Propagation}
\label{ch:ch14-constant-propagation}

Programmers write symbolic names where machines want concrete numbers.
Loop bounds arrive as configuration constants, feature flags arrive as command-line booleans, and table sizes arrive as header macros.
By the time control reaches the hot loop, many variables hold one fixed value along every path that matters.
The problem of this chapter is to discover those values flow-sensitively, substitute them into uses, and fold the operations they enable.
Discovery must be conservative: a variable becomes constant only when every reaching path agrees, and any disagreement demotes it permanently.
Section~\ref{sec:ch14-lattice} defines the value domain, Section~\ref{sec:ch14-algorithm} propagates it, and Section~\ref{sec:ch14-worked} computes the diamond in Figure~\ref{fig:ch14-diamond} by hand.

\section{A lattice of constant claims}
\label{sec:ch14-lattice}

A \term{constant lattice} for one variable holds top above every integer constant and every constant above bottom.
Top means \enquote{not yet known to be anything}, which seeds every variable before analysis.
A concrete constant such as $7$ means \enquote{known to be exactly this on all paths seen so far}.
Bottom means \enquote{known not to be any single constant}, which covers disagreement as well as genuinely varying values.
The meet of top with any value returns that value, because an unseen path contributes no evidence.
The meet of equal constants returns that constant, because all evidence agrees.
The meet of distinct constants returns bottom, because no one constant covers both.
Bottom absorbs every operand of meet.
An assignment can still establish a new constant at a later program point; the descending-chain statement applies to successive solver states at one fixed point.
Table~\ref{tab:ch14-meet} tabulates the meet over representative pairs, and the \enquote{once bottom} rule is the column to memorize.

Consider \code{d = 0} on the left arm and \code{d = load p} on the right arm.
The left contributes constant $0$, the right contributes bottom because a load is not a known constant.
Their meet is bottom, so the join block reports \code{d} as nonconstant.
Substituting zero at the join would miscompile every execution taking the right arm.
Now change the right arm to \code{d = 0} as well, and the meet of $0$ with $0$ stays $0$.
Substitution at the join is then safe, and folding \code{t = d + 1} to $1$ follows.
The lattice distinguishes these cases with one table lookup, which is why constant propagation scales where manual inspection does not.
The \enquote{load means bottom} default is deliberately harsh: memory may hold a constant today and another tomorrow, and only alias analysis can soften it.
We revisit that harshness in Section~\ref{sec:ch14-implementation} with SSA-based sparse variants.

\begin{irlisting}
; Diamond with agreeing and disagreeing arms for hand analysis
; B1 branches on flag, B2 and B3 assign d and e, B4 uses both
B1: br flag, B2, B3
B2: d = 7
B2: e = d + 1
B2: jmp B4
B3: d = 7
B3: e = d * 2
B3: jmp B4
B4: use d
B4: use e
\end{irlisting}

The listing sets up both outcomes in one graph: \code{d} agrees at $7$ on both arms, while \code{e} computes $8$ on the left and $14$ on the right.
A tally that treats every multiply-defined variable as varying loses the agreement for \code{d}; flow-insensitive analyses need not make that particular approximation.
A flow-sensitive propagation instead meets the arm values at $B4$ and keeps \code{d} constant while demoting \code{e}.
The uses in $B4$ then fold for \code{d} and stay symbolic for \code{e}.
No test run could establish this for all flags without enumerating the flag, yet the lattice does it in two passes.
Algorithm~\ref{alg:ch14-constprop} is the engine, and Figure~\ref{fig:ch14-diamond} draws the graph it traverses.

Transfer functions evaluate each operation in the lattice rather than in integers.
Addition of two constants folds to their sum, addition with bottom yields bottom, and addition with top waits by returning top.
A phi meets its arm values with the same table, which is why agreeing arms stay constant and disagreeing arms fall.
A load, call, or volatile read returns bottom unless a trusted domain fact says otherwise.
A store kills constant facts about aliased locations, conservatively demoting them to bottom.
Monotonicity holds throughout: refining an input from top to constant to bottom never refines an output backward.
Each scalar cell has height three; a state with $V$ cells has height $2V+1$, so the product still has finite height.
\section{Forward propagation over the lattice}
\label{sec:ch14-algorithm}

\begin{algorithm}
\caption{Forward constant propagation on the flat lattice}\label{alg:ch14-constprop}
\begin{pseudocode}
\Require entry-reachable CFG with predecessor-free entry, monotone transfers $\Tf{B}$, sound boundary $in_0$
\Ensure $\In{B}$ and $\Out{B}$ over-approximate runtime values; constants may be substituted
\State $\In{B} \gets \top$ and $\Out{B} \gets \top$ for every block $B$
\State $\In{entry} \gets in_0$
\State $W \gets$ all blocks in $G$
\While{$W \neq \emptyset$}
\State remove some $B$ from $W$
\State $in \gets \meet_{P \in \Pred(B)} \Out{P}$, or $in_0$ if $B$ is entry
\State $out \gets \Tf{B}(in)$
\If{$in \neq \In{B}$ or $out \neq \Out{B}$}
\State $\In{B} \gets in$
\State $\Out{B} \gets out$
\State add $\Succ(B)$ to $W$
\EndIf
\EndWhile
\end{pseudocode}
\end{algorithm}

Algorithm~\ref{alg:ch14-constprop} seeds every block with top and descends toward constants and bottom.
Seeding with top is the load-bearing choice: top contributes nothing to meets, so first evidence decides.
Seeding with bottom would instead poison every meet immediately and report nothing constant anywhere.
Each visit folds the block operations in lattice arithmetic and propagates the result along successors.
A cell moves at most twice, from top to constant to bottom; a block with many cells can need many visits.
Unknown runtime inputs must start at bottom, not top: lack of information about a parameter is not an unvisited path.
Require definite assignment before each use, or model the language's uninitialized-value semantics explicitly.
At a reachable use, a constant licenses substitution and bottom forbids it; top is a solver marker, not a runtime value that may be chosen arbitrarily.
Substitution rewrites uses of constant variables to the constant and deletes the defining copy when dead-code analysis agrees.
Folding then evaluates newly constant operations, which can turn branches into jumps and enable the cascade of the previous chapter.
Morgan organizes redundancy and motion passes around exactly this discovered-constant feed, see~\cite{morgan-1998-optimizing}.

Sparse conditional constant propagation, usually shortened to \term{SCCP}, improves on this dense form in two ways.
It tracks values per SSA name rather than per program point, which removes most block visits.
It also treats branches conditionally: an edge from a provably untaken arm contributes nothing to the meet at its join.
An over-eager dense pass meets both arms even when the condition is constant, demoting variables that only disagree along dead paths.
SCCP instead marks edges executable as constants emerge, and joins meet only over executable predecessors.
The result finds more constants at lower cost, at the price of interleaving branch folding with propagation.
We implement dense propagation first for clarity, then treat SCCP as the production upgrade in Section~\ref{sec:ch14-implementation}.
Appel presents the executable-edge idea with the data structures it needs, see~\cite{appel-1998-modern}.
\section{Worked example on the diamond}
\label{sec:ch14-worked}

We propagate by hand over Figure~\ref{fig:ch14-diamond} with the unknown input \code{flag} at bottom.
The not-yet-assigned locals \code{d} and \code{e} start at top; both arms define them before any use.
Block $B1$ passes these states through unchanged.
Block $B2$ evaluates \code{d} to $7$ and \code{e} to $8$ by folding $7 + 1$.
Block $B3$ evaluates \code{d} to $7$ and \code{e} to $14$ by folding $7 \times 2$.
Block $B4$ meets \code{d} as $7 \meet 7 = 7$ and \code{e} as $8 \meet 14 = \bot$.
Table~\ref{tab:ch14-meet} confirms each pair lookup, and Table~\ref{tab:ch14-trace} records the per-block lattice states across two passes.
The hand arithmetic below checks every meet and every fold explicitly.

\begin{figure}
\begin{tikzpicture}
\node[basic block] (b1) {B1: br flag};
\node[basic block, below left=1.1cm and 0.7cm of b1] (b2) {B2: d = 7, e = 8};
\node[basic block, below right=1.1cm and 0.7cm of b1] (b3) {B3: d = 7, e = 14};
\node[basic block, below=2.4cm of b1] (b4) {B4: use d, use e};
\draw[cfg edge] (b1) -- (b2);
\draw[cfg edge] (b1) -- (b3);
\draw[cfg edge] (b2) -- (b4);
\draw[cfg edge] (b3) -- (b4);
\end{tikzpicture}
\caption{The constant-propagation diamond where \code{d} agrees at $7$ on both arms and \code{e} disagrees, so only \code{d} folds at the join.}
\label{fig:ch14-diamond}
\end{figure}

\begin{table}
\caption{Meet table for the flat constant lattice over representative values, read as row meet column.}
\label{tab:ch14-meet}
\begin{tabular}{lccc}
\toprule
Meet & Top & Constant c & Bottom \\
\midrule
Top & top & constant c & bottom \\
Constant d & constant d & c if same else bottom & bottom \\
Bottom & bottom & bottom & bottom \\
\bottomrule
\end{tabular}
\end{table}

\begin{table}
\caption{Lattice states for \code{d} and \code{e} at each block exit across two propagation passes.}
\label{tab:ch14-trace}
\begin{tabular}{lcc}
\toprule
Block & Pass 1 exit values & Pass 2 exit values \\
\midrule
B1 & d top, e top & d top, e top \\
B2 & d 7, e 8 & d 7, e 8 \\
B3 & d 7, e 14 & d 7, e 14 \\
B4 & d 7, e bottom & d 7, e bottom \\
\bottomrule
\end{tabular}
\end{table}

Pass one visits $B1$ and holds top, visits $B2$ and folds to $(7, 8)$, visits $B3$ and folds to $(7, 14)$, then visits $B4$ and meets to $(7, \bot)$.
Pass two revisits every block: $B1$ unchanged, $B2$ recomputes $(7, 8)$ identically, $B3$ recomputes $(7, 14)$ identically, and $B4$ remeets to $(7, \bot)$ identically.
No value changes, so the worklist empties and the fixed point stands.
Substitution replaces \code{d} by $7$ at its uses in $B4$ and folds any operation on it, for example \code{d + 1} to $8$.
A \term{folding} step evaluates newly constant operations, which can turn branches into jumps and enable the cascade of the previous chapter.
It leaves \code{e} symbolic because bottom forbids substitution.
The \enquote{agreeing versus disagreeing} pattern is now a mechanical table lookup rather than an informal inspection.
A second example worth tracing changes $B3$ to compute \code{e} as $8$ as well, after which $B4$ keeps $(7, 8)$ and both uses fold.

The same graph exposes the dense-pass weakness that SCCP repairs.
Suppose $B1$ branches on a flag that earlier propagation proves true, so the edge to $B3$ never executes.
Dense propagation still meets $B3$ values into $B4$, demoting \code{e} to bottom even though no execution brings $14$ to the join.
SCCP instead marks only the true edge executable, meets over executable predecessors alone, and keeps \code{e} at $8$.
The difference matters whenever a flag guards a stale or exceptional arm with wild values.
Dense form is safe in both cases, but infeasible paths can reduce its precision even when its final constants remain correct.
Conditional form recovers the precision by refusing to pollute joins with dead arms.

\section{Implementation notes and cost}
\label{sec:ch14-implementation}

A dense implementation stores one lattice cell per variable per block, updated by the block transfer on each visit.
Let $N$ be blocks, $E$ edges, and $V$ variables.
Each cell changes at most twice, yielding $O(N+EV)$ worklist visits with duplicate-free queue membership.
A visit to $B$ additionally costs $O(V(|\Pred(B)|+1)+I_B)$ for dense meets and a scan of its $I_B$ bounded-arity instructions.
Multiply visit counts by these costs rather than treating a whole-state update as one cell update.
Memory is $O(N \cdot V)$ cells, which dominates for large functions and motivates the sparse upgrade.
SSA-based sparse propagation stores one cell per SSA name globally, cutting memory to $O(V)$ names and visits to definition-use edges only.
The executable-edge marking adds one flag per CFG edge with negligible cost.
Safonov discusses bounding such passes so trusted builds stay predictable, see~\cite{safonov-2010-trustworthy}.

Precision hazards cluster around memory, calls, and undefined behavior.
A store through a pointer demotes every potentially aliased variable to bottom unless alias analysis refines the may-alias set.
A call demotes globals and escaped locals for the same reason, unless interprocedural facts prove them preserved.
Signed overflow folds only when the language standard defines wrapping for the type; otherwise the folder must preserve the trap or the undefined latitude explicitly.
The C definition of sequence points, effective types, and trap representations governs each of these guards, see~\cite{wg14-2011-c11}.
% TODO(cite): need a primary reference for Wegman-Zadeck conditional techniques beyond the four listed books.

\section{Exercises}
\label{sec:ch14-exercises}

\begin{enumerate}
\item \label{ex:ch14-hand-variant} Recompute Table~\ref{tab:ch14-trace} on Figure~\ref{fig:ch14-diamond} after changing $B3$ to \code{d = 8}.
Show every meet with Table~\ref{tab:ch14-meet} and state which uses in $B4$ still fold.
Then change $B1$ to branch on constant true and explain why dense propagation still demotes \code{d} while SCCP keeps it.
Give the executable-edge set SCCP marks in each case.
\item \label{ex:ch14-bottom-poison} Seed the nonboundary states of Algorithm~\ref{alg:ch14-constprop} with bottom instead of top and run it on Figure~\ref{fig:ch14-diamond}.
Show why literal assignments still establish constants and why the diamond reaches the same result.
Then analyze a loop whose preheader assigns \code{x = 7} and whose body leaves \code{x} unchanged.
Show that the header equation $X=7\meet X$ admits both $7$ and bottom, and explain why the prescribed initialization selects the useful solution.
\item \label{ex:ch14-sparse-build} Implement dense propagation and an SSA-based sparse variant for the diamond family with $V$ variables and two arms.
Measure block visits and cell writes for $V$ set to $10$, $100$, and $1000$ with agreeing and disagreeing arms.
Report where sparse wins, where dense ties, and how executable-edge marking changes the disagreeing case.
Recommend which variant a just-in-time compiler should run under a tight time budget.
\end{enumerate}

\section{Further reading}
\label{sec:ch14-reading}

Appel implements the executable-edge marking that makes conditional propagation practical, see~\cite{appel-1998-modern}.
Morgan shows how discovered constants feed redundancy elimination and code motion, see~\cite{morgan-1998-optimizing}.
Safonov examines cost-bounded analysis design for trusted compilers, see~\cite{safonov-2010-trustworthy}.
The C language definition states the overflow, aliasing, and sequencing rules that folding guards must respect, see~\cite{wg14-2011-c11}.
